\begin{enumerate}[label=\textbf{\ansnum.\arabic*}]
  \item \texttt{w}+\texttt{q}: walk with $v_x=0.3$, $v_{\text{yaw}}=0.6$ (forward while turning left). \texttt{w}+\texttt{s}: $v_x=0.3-0.2=0.1$, a slow walk forward; a cleaner design ignores opposite keys held together.
  \item Commands are sent every $0.1$\,s. During a $0.25$\,s loss, the time between two \emph{received} commands is between $0.25$ and about $0.35$\,s, depending on where the gap falls between two sends. With a time to live of $0.3$\,s the robot usually keeps walking and at worst stops for a few hundredths of a second. With $0.15$\,s it stops at every such gap, for $0.1$ to $0.2$\,s: a stutter. A time to live of about three sending periods is a good compromise.
  \item
\begin{verbatim}
def test_drive_dropped_while_playing():
    slot = CommandSlot()
    server = DebugServer(slot, port=0)
    server.game_state = "playing"
    server.handle(pb.Request(command=pb.MotionCommand(
        kind="walk", vx=0.3, ttl=0.3, priority=10)))
    assert slot.take(time.monotonic()) is STAND
\end{verbatim}
(This needs the server's request handling in a method \texttt{handle} that a test can call directly, and a \texttt{game\_state} attribute that cognition updates.)
  \item A stop that lasts for ever would need a separate ``release'' command, and if that is forgotten, the robot never moves again and nobody knows why. Five seconds is long enough to walk to the robot; pressing \texttt{x} again renews it.
  \item
\begin{verbatim}
def command_for(keys_down, speed=1.0):
    if "shift" in keys_down:
        speed *= 0.5
    ...                         # as before
\end{verbatim}
\end{enumerate}
